\documentclass[a4paper, 11pt]{article}

\usepackage[a4paper, top=2.5cm, bottom=2.5cm, left=2cm, right=2cm]{geometry}
\usepackage{graphicx} 
\usepackage[T5]{fontenc}
\usepackage[utf8]{inputenc}
\usepackage[vietnamese, english]{babel}
\usepackage{helvet}

\renewcommand{\familydefault}{\sfdefault}

\title{\textbf{Brief Report: TCP/IP File Transfer}}
\author{Nguyễn Mai Hùng}
\date{\today}

\begin{document}
\maketitle

\section{Overview}
Practical work 1 which require developing a CLI-based tool to transfer files via TCP/IP. The system uses a Client-Server model where the Client streams binary data and the Server listens on port 5000 to reconstruct the file.

\section{Architecture}
\begin{itemize}
    \item \textbf{Server:} Binds to 0.0.0.0:5000, accepts connections, performs a handshake (filename/size), and writes binary chunks to disk.
    \item \textbf{Client:} Connects to the server, sends metadata, and streams file content.
\end{itemize}

\section{Troubleshooting}
During testing, files were successfully sent but failed to appear in the server's directory.
\begin{itemize}
    \item \textbf{Issue:} The server saved files relative to the terminal's working directory, often misplacing them in the root user folder.
    \item \textbf{Solution:} Implementation of a absolute path handling to enforce saving files in the script's exact directory.
\end{itemize}

\section{Conclusion}
The application is now stable. The absolute path fix ensures files are saved reliably regardless of how the server script is launched.

\end{document}